\subsection{Prefix restrictions and multilinear evaluation}\label{sec:kbfoldfold:prefix}

Throughout this chapter, $\F$ is a field, $m\geq1$ is an integer, $N=2^m$, and $I_k=\{0,1\}^k$ for every $0\leq k\leq m$. The first bit is the lowest-index bit, as in Definition~\ref{def:intro:indices}. The variables of a table are ordered, and no permutation of their indices is implicit.

\begin{definition}[Prefix equality weight]\label{def:kbff:prefix-weight}
For $0\leq j\leq m$, $a\in I_j$ and $r=(r_0,\ldots,r_{j-1})\in\F^j$, define
\[
 \lambda^{(j)}_a(r)
 :=\prod_{i=0}^{j-1}\bigl((1-a_i)(1-r_i)+a_ir_i\bigr)\in\F.
\]
For $j=0$, the unique weight is the empty product $1$.
\end{definition}

\begin{lemma}[First-bit factorization]\label{lem:kbff:weight-factor}
For $j<m$, any $a=(a_0,\ldots,a_{j-1})\in I_j$, $e\in\{0,1\}$ and $r'=(r_0,\ldots,r_{j-1})\in\F^j$,
\[
 \lambda^{(j+1)}_{(a,e)}(r',t)=
 \lambda^{(j)}_a(r')\bigl((1-e)(1-t)+et\bigr).
\]
\end{lemma}
\begin{proof}
By Definition~\ref{def:kbff:prefix-weight}, the left side is a product indexed by $i=0,\ldots,j$. Separate the factor with $i=j$ from the product with $i<j$. The latter is $\lambda^{(j)}_a(r')$ and the former is the displayed affine factor.
\end{proof}

\begin{definition}[Simultaneous prefix restriction]\label{def:kbff:prefix-restriction}
For $0\leq j\leq m$, $r\in\F^j$, and $f:I_m\to\F$, define
\[
 \mathsf{Res}^{[j]}_r(f):I_{m-j}\to\F,\qquad
 \mathsf{Res}^{[j]}_r(f)(b)
   :=\sum_{a\in I_j}\lambda^{(j)}_a(r)\,f(a,b).
\]
Here $(a,b)\in I_m$ denotes concatenation in the original bit order.
\end{definition}

\begin{lemma}[Zero-prefix identity]\label{lem:kbff:prefix-zero}
For every table $f:I_m\to\F$,
\[
 \mathsf{Res}^{[0]}_{()}(f)=f.
\]
\end{lemma}
\begin{proof}
The set $I_0$ contains only the empty tuple $()$, whose equality weight equals $1$. In the defining sum, $((),b)=b$. Thus $\mathsf{Res}^{[0]}_{()}(f)(b)=f(b)$ for every $b$.
\end{proof}

\begin{lemma}[Prefix extension by one restriction]\label{lem:kbff:prefix-step}
For $0\leq j<m$, $r\in\F^j$, and $t\in\F$,
\[
 \mathsf{Res}^{[j+1]}_{(r,t)}(f)
 =\Res_t\bigl(\mathsf{Res}^{[j]}_r(f)\bigr).
\]
\end{lemma}
\begin{proof}
Fix $b\in I_{m-j-1}$ and expand the left side:
\[
 \mathsf{Res}^{[j+1]}_{(r,t)}(f)(b)
 =\sum_{a\in I_j}\sum_{e\in\{0,1\}}
   \lambda^{(j+1)}_{(a,e)}(r,t) f(a,e,b).
\]
By Lemma~\ref{lem:kbff:weight-factor}, the summand equals
$\lambda^{(j)}_a(r)((1-e)(1-t)+et)f(a,e,b)$.
Separate the two values of $e$:
\[
 (1-t)\sum_{a\in I_j}\lambda^{(j)}_a(r)f(a,0,b)
 +t\sum_{a\in I_j}\lambda^{(j)}_a(r)f(a,1,b).
\]
The two sums are $\mathsf{Res}^{[j]}_r(f)(0,b)$ and
$\mathsf{Res}^{[j]}_r(f)(1,b)$, respectively. Their affine combination is
$\Res_t(\mathsf{Res}^{[j]}_r(f))(b)$ by Definition~\ref{def:intro:restriction}.
Since $b$ was arbitrary, the table identity follows.
\end{proof}

\begin{lemma}[Linearity of prefix restriction]\label{lem:kbff:prefix-linear}
For any $f,g:I_m\to\F$, $c\in\F$, and every prefix $r\in\F^j$,
\[
 \mathsf{Res}^{[j]}_r(f+cg)=
 \mathsf{Res}^{[j]}_r(f)+c\,\mathsf{Res}^{[j]}_r(g).
\]
\end{lemma}
\begin{proof}
At each remaining bit tuple $b$, distribute the sum
$\sum_a\lambda^{(j)}_a(r)(f(a,b)+cg(a,b))$, and take the scalar $c$ outside the finite sum.
\end{proof}

\begin{lemma}[Partial multilinear-evaluation identity]\label{lem:kbff:prefix-mle}
For every $0\leq j\leq m$, $r\in\F^j$, and $y\in\F^{m-j}$,
\[
 \widetilde{\mathsf{Res}^{[j]}_r(f)}(y)=\widetilde f(r,y).
\]
\end{lemma}
\begin{proof}
By the definition of multilinear extension, the left side is
\[
 \sum_{b\in I_{m-j}}
     \mathsf{Res}^{[j]}_r(f)(b)\lambda^{(m-j)}_b(y).
\]
Substitute Definition~\ref{def:kbff:prefix-restriction} and interchange the two finite sums:
\[
 \sum_{a\in I_j}\sum_{b\in I_{m-j}}
 f(a,b)\lambda^{(j)}_a(r)\lambda^{(m-j)}_b(y).
\]
By splitting the product in Definition~\ref{def:intro:mle},
$\lambda^{(m)}_{(a,b)}(r,y)=
\lambda^{(j)}_a(r)\lambda^{(m-j)}_b(y)$.
The map $(a,b)\mapsto(a,b)$ is a bijection
$I_j\times I_{m-j}\to I_m$; thus the displayed expression equals
$\sum_{c\in I_m}f(c)\lambda^{(m)}_c(r,y)=\widetilde f(r,y)$.
\end{proof}

\begin{corollary}[Complete prefix restriction]\label{cor:kbff:full-prefix}
For $r\in\F^m$ the table $\mathsf{Res}^{[m]}_r(f)$ has one entry and
\[
 \mathsf{Res}^{[m]}_r(f)(())=\widetilde f(r).
\]
\end{corollary}
\begin{proof}
Apply Lemma~\ref{lem:kbff:prefix-mle} with $j=m$ and the unique empty tuple of $\F^0$; the multilinear extension of a one-entry table is its only entry.
\end{proof}
